% UBT-AI-PROVENANCE-BEGIN
% schema: ubt-ai-provenance/v1
% tier: C_working
% ai_assistance: disclosed
% human_review: risk-based
% editorial_responsibility: Ing. David Jaroš
% policy: ../../../AI_PROVENANCE.md
% notice: Working material; exhaustive human review is not claimed.
% UBT-AI-PROVENANCE-END

\chapter{Program konstanty jemné struktury: co je odvozeno a co není}
\label{ch:alpha}

\section{Status před algebrou}
Autoritativní stav úložiště je záměrně silnější než jakýkoli elegantní
vzorec v této kapitole:
\begin{quote}
Fyzikální konstanta jemné struktury ještě není odvozena z prvních principů
v UBT.  Cesta integer-$137$ je strukturální/podmíněná, zatímco most
stanovení požadovaného efektivního koeficientu z akce UBT zůstává otevřené.
\end{quote}
Živé zdroje stavu jsou \texttt{STATUS.md},
\texttt{WHAT\_IS\_PROVED.md} a
\texttt{canonical/alpha/ALPHA\_MASTER\_STATUS.md}. Historické texty používající
fráze ``odvození alfa bez přizpůsobení'' nepřepisují tyto účetní knihy.

\section{Podmíněný winding potenciál}
Centrální redukované modelové studie
\[
 \boxed{V_{\rm eff}(n)=n^2-Bn\ln n,\qquad n>0.}
\]
Dočasné zacházení s $n$ jako s kontinuální proměnnou,
\[
 \frac{dV_{\rm eff}}{dn}
 =2n-B(\ln n+1).
\]
Stacionární bod tedy vyhovuje
\[
 \boxed{B=\frac{2n}{\ln n+1}.}
\]
Pokud někdo požádá o stacionární bod na $n=137$, požadovaný koeficient je
\[
 B_{137}=\frac{274}{\ln137+1}\approx46.284.
\]
Tato rovnice je přesnou vlastností redukovaného potenciálu. \emph{Nevysvětluje},
proč musí akce UBT generovat právě tuto hodnotu $B$.

Druhá derivace je
\[
 \frac{d^2V_{\rm eff}}{dn^2}=2-\frac{B}{n}.
\]
Ve stacionárním bodě se to stane
\[
 2-\frac{2}{\ln n+1}
 =\frac{2\ln n}{\ln n+1}>0
 \qquad(n>1),
\]
takže spojitý stacionární bod je lokální minimum.  Celé číslo a
omezení primárního sektoru pak lze studovat jako další diskrétní
problém s výběrem.

\section{Co může a nemůže znamenat ``prvočíselná stabilita''}
Úložiště obsahuje strukturální analýzu primární stability, ve které
celočíselný sektor kolem $137$ je zkoumán podle specifikovaného pravidla výběru.  A
prvotřídně označené minimum může být reprodukovatelným výsledkem tohoto zmenšeného modelu.
Však:
\begin{itemize}
 \item Základem $137$ je standardní teorie čísel;
 \item výběr primárního sektoru sám o sobě neidentifikuje
       $\alpha^{-1}$ s tímto sektorem;
 \item fyzická oprava z $137$ na $137.036\ldots$ vyžaduje samostatnou
       derivace;
 \item Koeficient $B$ je stále třeba získat spíše z dynamiky UBT
       než je odvozeno z cílové hodnoty.
\end{itemize}
Kapitola theta/Mellin/Feynman rozvíjí možné mechanismy porozumění
kam by mohla vstoupit multiplikativní primární struktura, ale tyto mechanismy jsou
výzkumné mosty, nedokončené alfa důkazy.

\section{Audit \(N_{\rm eff}\): dvě různá počítání se nesmějí zaměňovat}
Použitá předchozí derivace
\[
 N_{\rm eff}=3\times2\times2=12.
\]
Současný audit rozlišuje alespoň dva objekty:
\begin{align*}
 N_{\rm eff}^{\rm twist}&=12,
 &&\text{in the SU(2)-twist construction},\\
 N_{\rm eff}^{\rm loop}&=3,
 &&\text{in the direct audited scalar-action loop count}.
\end{align*}
Identifikace
\[
 N_{\rm eff}^{\rm loop}=N_{\rm eff}^{\rm twist}
\]
není v současné době dokázaná věta.  Proto vzorec, který vkládá $12$ do
koeficient smyčky musí pojmenovat předpoklady zkroucení; nelze to popsat jako
jedinečný důsledek samotného $\mathbb C\otimes\mathbb H$.

V souladu s tím úložiště zaznamenává dva různé koeficienty základní linie
trasy,
\[
 B_0^{\rm loop}=2\pi,
 \qquad
 B_0^{\rm twist}=8\pi,
\]
s můstkem z obou základních linií k fenomenologicky požadovanému
$B\approx46.284$ stále otevřený.

\section{Mezera G137-B}
Rozhodující alfa problém lze říci bez rétoriky:
\begin{quote}
Vycházejte z kanonické akce UBT a jejího oprávněného obsahu pole
efektivní redukovaný koeficient násobící $n\ln n$, včetně všech
normalizace a multiplicity a ukazují, že odvození je nezávislé
naměřené hodnoty $\alpha$.
\end{quote}
Toto je Gap G137-B.  Vzniklo několik historických a výzkumných tras
zajímavé modulární, eta, gama, determinantové a vinuté struktury, ale
aktuální master status zaznamenává požadované vložení do fyzického
koeficient jako otevřený.

\section{Proč má stále smysl studovat theta/Mellinovu větev}
Selhání naivního řetězce
\[
 \Theta\to\text{Mellin}\to\zeta\to\text{Euler product}
 \to\text{prime sectors}
\]
do\emph{dynamicky} vybrat prvočísla je samo o sobě užitečné.  Produkt Euler společnosti a
standardní theta Mellinova transformace již odráží multiplikativní aritmetiku
přítomný v sadě celočíselných indexů.  Musel by se identifikovat silnější mechanismus
primitivní sektory dynamicky.

Racionální-časové theta oživení jsou jedním z kandidátů, protože kvadratické Gaussovy součty
faktor nad coprime jmenovatele.  V tomto nastavení jsou hlavní síly skutečně
neredukovatelné stavební kameny aritmetiky obrození.  Ať už alfa
konstrukce primárního sektoru může být přepsána v těchto proměnných je explicitní
otevřená otázka, nikoli stanovený výsledek.

\section{Bezpečné shrnutí}
Výroky k zapamatování jsou:
\begin{enumerate}
 \item redukovaná stacionární rovnice a její výpočet stability jsou přesné;
 \item strukturní výběr celé číslo/prvočíslo lze studovat reprodukovatelně;
 \item $B$ ještě není odvozen z akce UBT na požadované hodnotě;
 \item příslušné multiplicitní účetnictví není plně odsouhlaseno;
 \item $\alpha^{-1}=137.036\ldots$ dnes není výsledkem prvního principu UBT.
\end{enumerate}
To je záměrně méně dramatické než starší archivní text a užitečnější pro
budoucí práci, protože nám přesně říká, která rovnice musí být uzavřena jako další.
